\chapter{Introduction}

% Background on 2025 Sumatra flash flood, introduce extreme event attribution study
% In late November 2025, Sumatra was devastated by exceptional monsoon rains and a rare tropical storm by Tropical Cyclone Senyar (Figure \ref{fig:domain}) with lifetime between 25 and 28 November 2025. The extreme rainfall triggered massive flooding, flash floods, and widespread landslides across provinces, including West Sumatra, North Sumatra and Aceh, with more than 3 million people have been affected. Given the high impact event, we might wonder the extent of climate change role in driving this event. Extreme event attribution study is an emerging field that equips us with the methods to explore this question.
At the end of November 2025, Sumatra was devastated by exceptional monsoon rains and a rare equatorial tropical storm, Tropical Cyclone (TC) Senyar \citep{climameter,egusphere-2026-694}. Forming near the Strait of Malacca—a tropical basin long considered almost cyclone-free due to the weak Coriolis force at low latitudes—Senyar’s development was enabled by a strong monsoon surge and exceptionally active convective clustering \citep{yulihastin2026origins}. The storm moved towards Sumatra, allowing convective activities to persist over land and triggering massive flooding, flash floods, and widespread landslides that affected more than 3 million people \citep{wu2026analysis, wwa, climameter}, as summarized in Figure \ref{fig:impact_marker}.

% Summarize past rapid attribution study to the same event
% Rapid attribution studies have assessed the likelihood of anthropogenic climate change (ACC) influencing the event. Based on on a ClimaMeter study that utilizes analogue-based approach, this event was associated with exceptional meteorological conditions and can be ascribed to ACC \citep{climameter}. Furthermore, meteorological conditions similar to those causing the Indonesian floods are up to 10\% wetter in the present, compared to the past. Another rapid study conducted by World Weather Attribution, which utilizes probabilistic approach, detected increase in extreme rainfall associated with rising global mean surface temperature (GMST) is estimated to be 9\%--50\% \citep{wwa}. However, the study could not attribute this trend to ACC as climate models do not show consistent results of change directions. This indicated atmospheric circulation might be systematically misrepresented by the models.
Rapid attribution studies have already assessed the likelihood of anthropogenic climate change (ACC) influencing the event. A study by ClimaMeter using an analogue-based approach suggested that meteorological conditions causing the Sumatra floods are up to 10\% wetter in the present compared to the past \citep{climameter}. Furthermore, a probabilistic assessment by World Weather Attribution (WWA) detected an increase in extreme rainfall associated with historical warming estimated at 9\%–50\% \citep{wwa}. However, WWA could not robustly attribute this trend to ACC because climate models showed inconsistent results regarding the direction of change. This indicated that atmospheric circulation and tropical convective processes might be systematically misrepresented by coarse-resolution global models.

% Introduce physical-based attribution study and relevant past studies as benchmark (2024 Valencia paper, Major tropical cyclone events paper)
% While these rapid attribution studies allow quick communication to public and quantification of the event likelihood change, there are limitation to this unconditional approaches which treats an extreme event as part of a class of similar events. Individual events can deviate from average behavior and their impacts often depend on local, nonlinear processes that are not captured when events are grouped into classes \citep{van2026essential}. Conditional approaches, also known as storyline methods, offer another perspective by fixing the large-scale meteorological condition and see how the event evolves in a warmer or cooler climate. While this physical-based methods require more time to be conducted, it provides insights that cannot be offered through the unconditional approaches. In performing conditional attribution study, pseudo-global warming (PGW) approach has been commonly used \citep{brogli2023pseudo, calvo2026human}, including for tropical cyclones \citep{patricola2018anthropogenic}.
While these rapid attribution studies allow for quick communication to the public, there are limitations to these "unconditional" approaches, which treat an extreme event as part of a broad class of similar events \citep{lenderink2025pseudo,van2026essential}. Individual events can deviate from average behavior, and their impacts often depend on local, nonlinear processes and specific mesoscale dynamics that are not captured when events are grouped into classes \citep{van2026essential}. Conditional approaches, also known as storyline methods, offer another perspective by fixing the large-scale meteorological conditions to see how a specific event evolves in a warmer or cooler climate.

% With pseudo-global warming method, perturbations are applied in the boundary conditions of a simulation, such that an event is simulated at different climate. Similar strategy was used in \citet{lenderink2025pseudo} to study convective extremes on past colder and future warmer climate conditions. By this, beyond attribution study, we can study the future projection of the event.
In performing conditional attribution studies, the pseudo-global warming (PGW) approach has been commonly used \citep{brogli2023pseudo, calvo2026human}, including for studies on tropical cyclones \citep{patricola2018anthropogenic}. With the PGW method, thermodynamic perturbations are applied to the boundary conditions of a simulation, allowing an event to be simulated across different climate states. Similar strategies have been used to study convective extremes in both past cooler and future warmer climates, allowing us to move beyond historical attribution toward future projection \citep{lenderink2025pseudo}.

Furthermore, while rainfall is the primary driver of Sumatra's flash floods, a comprehensive assessment must also consider changes in TC Senyar intensity. Previous storyline studies on major tropical cyclones have demonstrated that while ACC significantly enhances rainfall, the response of maximum surface wind speeds and minimum sea level pressure can vary depending on the storm’s environment and structure \citep{patricola2018anthropogenic}. By assessing TC Senyar across a range of climates—from the pre-industrial past to the current climate and a future +1.5$^\circ$ warming level—this study aims to determine how both the hydrological impact and the strength of equatorial TC Senyar evolve with global warming.

% Define research questions
By this, we define the following research objectives to attribute ACC to the 2025 Sumatra's flood-inducing extreme rainfalls and its future projection:
\begin{enumerate}
    \item \textbf{Event Reconstruction:} Simulate the 2025 extreme rainfall and cyclone dynamics associated with TC Senyar in Sumatra with high resolution RegCM5 model under present climate. The simulation shall be checked whether it can represent the observed storm or extreme precipitation well.
    \item \textbf{Historical Attribution:} Simulate a counterfactual model where the warming signal is removed from the boundary conditions to represent a pre-industrial climate.
    \item \textbf{Future Risk Projection:} Simulate a future projection by adding a +1.5$^\circ$ warming signal to the boundary conditions to assess the hazard's evolution in a warmer world.
    \item \textbf{Physical Mechanism Identification:} Quantify the the thermodynamic (e.g CAPE, evaporation) and dynamical (e.g updraft, moisture transport) changes between the models. 
\end{enumerate}

In general, under the PGW approach, we are asking: \textit{"If an event like the 2025 Sumatra floods had occurred under a pre-industrial (cooler) or future (warmer) climate, how would it have changed?"}

\begin{figure}[h]
\centering
\includegraphics[width=0.6\linewidth]{../../figs/event_view.png}
\caption{Accumulated rainfall during TC Senyar lifetime from MSWEP and the TC track from IBTrACS with impact markers}
\label{fig:impact_marker}.
\end{figure}